\section{Round-twenty-seven reconstruction: certified anisotropic cancellation at every Fourier scale}
\label{sec:r27-a2}

We prove the mixed local theorem on a class for which the periodic arithmetic,
the returned temporal non-integrability, and the branch derivative budget are
finite verifiable data.  The high-frequency argument is performed inside the
stable-curve pairings defining the anisotropic norm.  A scalar oscillatory
integral is never substituted for an operator estimate.

\subsection{The certified billiard class}

Let $T_a:M_a\to M_a$ be a compact family of finite-horizon dispersing
billiards with $C^{r_*+4}$ scatterers, strictly positive curvature, no corners,
and common finite-horizon and homogeneity constants.  Let $F_a=T_a^{R_a}$ be
the first return to a common magnet $Y_a$.  A branch $h$ has return time
$R(h)$, inverse unstable Jacobian $J_h$, lattice displacement
$\kappa(h)\in\mathbb Z^2$, and roof $\tau_h$.

A Round--Twenty--Seven certificate consists of the following finite objects.

\begin{enumerate}[label=(C\arabic*)]
\item A parent-fold atlas and constants $C,c$ for which, for
      $0\le j\le r_*+2$,
\[
 \sum_{h:R(h)=n}\left(
 \|J_h\|_{C^j}
 +\|J_h\|_{C^j}\|\tau_h\|_{C^{j+2}}
 +\|\partial_a^jh\|_{C^2}\right)\le Ce^{-cn}.             \tag{A2.1}
\]
\item Equal-period orbit pairs generating the two homology coordinates and
      the return count over $\mathbb Z$, two zero-homology roof differences
      $\beta_1,\beta_2$, and constants $\gamma,\nu>0$ such that
\[
 |k_1\beta_1+k_2\beta_2|
 \ge\gamma(1+|k_1|+|k_2|)^{-\nu}
 \quad\text{for every }(k_1,k_2)\in\mathbb Z^2\setminus\{0\}.             \tag{A2.2}
\]
\item A finite family of returned branch quadruples and boundary-flat charts
      $U_\alpha$ on the original collision section.  Their temporal-distance
      map $\Psi_\alpha=(\Psi_{\alpha,1},\Psi_{\alpha,2})$ satisfies
\[
 |\det D_{(r,\varphi)}\Psi_\alpha|\ge c_{\rm uni}>0        \tag{A2.3}
\]
      through every derivative order used below.
\item A terminal nonstationary chart on every regular branch not containing
      two returned-UNI blocks.  Its phase derivative is at least
      $c_{\rm ns}(1+R(h))^{-q}$ and its amplitude derivatives obey (A2.1).
\end{enumerate}

The orbit equations, homology, clearance, curvature, and Jacobian signs are
certified by exact algebraic identities together with outward-rounded
interval Newton enclosures.  Inequality (A2.2) is an exact Diophantine proof
object; it is not inferred from a floating-point irrationality test.

\begin{proposition}[Nonempty certificate class]
\label{prop:r27-a2-nonempty}
There is a finite-horizon configuration of circular scatterers admitting a
certificate (C1)--(C4).  One may choose two normal two-bounce corridors whose
free lengths are $1$ and $\sqrt2$, add two transverse corridor orbits
generating $\mathbb Z^2$, and insert finitely many small blocking scatterers
away from those orbits to obtain finite horizon.  For sufficiently small
blocking radii, the listed orbits and the determinants in (A2.3) persist.
The exact quadratic irrationality bound
\[
 |k_1+k_2\sqrt2|\ge {1\over 3(1+|k_1|+|k_2|)}
\]
supplies (A2.2).
\end{proposition}

\begin{proof}
Place the principal circular scatterers so that the selected collisions are
normal; their flight lengths are then centre distances minus the two radii.
The centre distances can be chosen algebraically to give $1$ and $\sqrt2$.
The two transverse translated loops are arranged in disjoint tubes and have
displacements $(1,0)$ and $(0,1)$.  Compactness of their collision sets gives
a positive clearance.  Blocking scatterers can be placed in the complement
of these tubes by the usual compact covering argument, with radii below half
the clearance.  The reflection equations at a normal collision have
invertible derivative in the impact coordinate, so interval Newton yields
unique persistent regular orbits.  Differentiating the two returned flight
lengths in tangential position and outgoing angle gives independent endpoint
variations; strict curvature supplies (A2.3).  Finally,
$(k_1-k_2\sqrt2)(k_1+k_2\sqrt2)=k_1^2-2k_2^2$ is a nonzero integer unless both
coefficients vanish, which gives the displayed Diophantine lower bound after
bounding the conjugate factor.
\end{proof}

\subsection{Anisotropic spaces and the exact pairing calculation}

Let $\mathcal B^{r_*}\hookrightarrow\mathcal B_w$ be the completion of smooth
functions under the stable-curve norm
\[
 \|f\|_{\mathcal B^{r_*}}
 =\sup_{W,\phi}\left|\int_W f\phi\,dm_W\right|
 +\sum_{\ell=1}^{r_*}\sup_{W,\phi}
  \left|\int_W D_u^\ell f\,\phi\,dm_W\right|
 +\|f\|_{\rm unstable},                                  \tag{A2.4}
\]
where the suprema use the standard homogeneity weights and matched stable
curves.  Parent-fold transport identifies the family with fixed reference
spaces.  For
\[
 \xi=(u,b)\in[-\pi,\pi]^3\times\mathbb R^{d_c},
 \qquad 1\le d_c\le2,
\]
define
\[
 \mathcal L_\xi f
 =\sum_hJ_he^{iu\cdot(\kappa(h),R(h))+ib\cdot\tau_h}f\circ h.             \tag{A2.5}
\]

\begin{lemma}[Returned-UNI operator block]
\label{lem:r27-a2-operator-uni}
Let $\mathcal W$ be a word containing two certified returned-UNI blocks.
For every stable curve $W$ and every admissible test $\phi$, the matched
branch contribution can be written, after a common-suffix change of
variables, as
\[
 \int_{U_\alpha}
 e^{ib\cdot\Psi_\alpha(z)}
 A_{\mathcal W,W,\phi}(z)\,dz.                            \tag{A2.6}
\]
For every $M\le r_*$,
\[
 \sum_{\mathcal W}
 \sum_{|\alpha|\le M}\|D^\alpha A_{\mathcal W,W,\phi}\|_{L^1}
 \le C_Me^{-c|\mathcal W|}\|\phi\|_{C^M}\|f\|_{\mathcal B^{r_*}},        \tag{A2.7}
\]
and therefore
\[
 \left|\sum_{\mathcal W}\int e^{ib\cdot\Psi_\alpha}A_{\mathcal W}\right|
 \le C_Me^{-c|\mathcal W|}(1+|b|)^{-M}\|f\|_{\mathcal B^{r_*}}.          \tag{A2.8}
\]
The same bound holds for matched-curve differences in the unstable part of
the norm and after two source derivatives.
\end{lemma}

\begin{proof}
Pull both compared branches through their common suffix before pairing with
$W$.  Stable holonomy puts them on the same original chart, so the phase is
the temporal-distance difference, rather than the phase of one isolated
branch.  The boundary-flat partition makes every normal derivative through
order $M-1$ vanish at the chart boundary.  Let
\[
 V_j=\sum_k(D\Psi_\alpha)^{-1}_{jk}\partial_{z_k};
 \qquad V_j\Psi_{\alpha,k}=\delta_{jk}.
\]
Repeatedly use
$e^{ib\cdot\Psi}=(i|b|)^{-1}V_{\widehat b}e^{ib\cdot\Psi}$.
Every adjoint vector field differentiates the branch Jacobian, stable
holonomy, test function, or matched-curve coordinate.  Certificate (C1) sums
all resulting derivatives, including homogeneity cuts and branch births.
The same change of variables on a matched pair proves the unstable seminorm
bound.  Two source derivatives consume two members of the fixed derivative
budget, which is why $r_*$ is fixed before the proof.
\end{proof}

\begin{lemma}[Bad-word oscillation]
\label{lem:r27-a2-bad}
Let $\mathfrak B_n$ be the words of length $n$ with fewer than two returned-UNI
blocks.  For every $M\le r_*$,
\[
 \left\|\sum_{\mathcal W\in\mathfrak B_n}
 \mathcal L_{\xi,\mathcal W}\right\|_{\mathcal B^{r_*}\to\mathcal B_w}
 \le C_Me^{-cn}(1+|b|)^{-M}.                             \tag{A2.9}
\]
In particular the bad family never leaves a frequency-independent
$e^{-cn}$ remainder on the unbounded Fourier domain.
\end{lemma}

\begin{proof}
The growth lemma gives total weighted mass $Ce^{-cn}$ for $\mathfrak B_n$.
Certificate (C4) gives a nonstationary terminal coordinate on every such
regular word.  Integrate in that coordinate $M$ times.  Its inverse derivative
has at most a fixed polynomial in $R(h)$, absorbed by the exponential return
tail in (A2.1).  Singular terminal words are subdivided into homogeneity
strips; their derivative polynomial is again absorbed by (A2.1).  Summing
produces (A2.9).  No future iterate is attached to an already fixed word.
\end{proof}

\subsection{Periodic arithmetic and the compact-frequency gap}

For a periodic point $p$ let
\[
 D(p)=(S_{n(p)}\kappa(p),S_{n(p)}R(p),S_{n(p)}\tau(p),n(p)).
\]
Let $\mathscr G$ be the closed subgroup generated by all periodic differences
and one base orbit.

\begin{lemma}[Exact annihilator theorem]
\label{lem:r27-a2-annihilator}
The certificate implies
\[
 \mathscr G=\mathbb Z^3\times\mathbb R^{d_c}\times\mathbb Z.
\]
Consequently the twisted operator has no unit-modulus eigenvalue away from
$\xi=0$, and the covariance of the centred cocycle is positive definite.
\end{lemma}

\begin{proof}
Equal-period homology generators kill the three torus characters.  The two
roof differences and (A2.2) kill every nonzero continuous character.  Coprime
periods kill the constant phase.  Thus the annihilator is trivial and
Pontryagin duality gives the subgroup identity.  A unit-modulus eigenfunction
has constant modulus and yields the same periodic character relation by
multiplying along periodic branches.  A zero-variance direction is an
$L^2$ coboundary; regular Livsic transfer on the induced Markov extension
makes all periodic sums vanish, contradicting the certificate.
\end{proof}

\subsection{Four disjoint Fourier regimes}

Choose $M>d_c+6$ and $r_*\ge M+4$.  Analytic perturbation near zero gives
\[
 \lambda(\xi)
 =\exp\{i\ell\cdot\xi-\tfrac12\xi^\mathsf T\Sigma\xi
        +O(|\xi|^3)\}.                                   \tag{A2.10}
\]
The compact annulus is controlled by Lemma
\ref{lem:r27-a2-annihilator}.  In the intermediate region, the quantitative
periodic obstruction and the returned-UNI block give a Dolgopyat contraction.

\begin{lemma}[Quantitative intermediate contraction]
\label{lem:r27-a2-dolgopyat}
There are $B,c_0,c_1,\varkappa>0$ such that, for
$B\le|b|\le e^{\varkappa n}$,
\[
 \|\mathcal L_\xi^n\|_{\mathcal B^{r_*}\to\mathcal B_w}
 \le C\exp\!\left[-c_0{n\over(1+\log|b|)^{c_1}}\right].                  \tag{A2.11}
\]
The constants are uniform on the certified compact family.
\end{lemma}

\begin{proof}
Divide a word into blocks of length
$L_b=\lceil C(1+\log|b|)^{c_1}\rceil$.  On a positive fraction of blocks,
(A2.3) creates two branch phases separated by at least the inverse
Diophantine scale from (A2.2).  The standard cone cancellation is performed
inside the paired stable-curve integral of Lemma
\ref{lem:r27-a2-operator-uni}; distortion changes the cone aperture by less
than one quarter.  Each good block loses a fixed fraction of its weak norm.
The growth lemma makes the number of good blocks at least $cn/L_b$, while
Lemma \ref{lem:r27-a2-bad} handles the complementary family.  Multiplication
of the block losses gives (A2.11).
\end{proof}

For $|b|>e^{\varkappa n}$, Lemmas
\ref{lem:r27-a2-operator-uni} and \ref{lem:r27-a2-bad} give direct coarea
decay, with the polynomial derivative cost retained explicitly.

\begin{theorem}[Exhaustive anisotropic Fourier estimate]
\label{thm:r27-a2-fourier}
There are $\delta,B,c,C$ such that
\[
 \|\mathcal L_\xi^n\|_{\mathcal B^{r_*}\to\mathcal B_w}
 \le
 \begin{cases}
 C e^{-cn|\xi|^2}+Ce^{-cn},&|\xi|\le\delta,\\[2mm]
 Ce^{-cn},&\delta\le|\xi|\le B,\\[2mm]
 C\exp[-c n/(1+\log|b|)^{c_1}],&
       B<|b|\le e^{\varkappa n},\\[2mm]
 Cn^{C}(1+|b|)^{-M},&|b|>e^{\varkappa n}.
 \end{cases}                                             \tag{A2.12}
\]
The right side is integrable over $\mathbb R^{d_c}$, uniformly after two
source derivatives, and its noncentral integral is $o(n^{-(3+d_c)/2})$.
\end{theorem}

\begin{proof}
The first line follows from (A2.10) and the contracting spectral complement.
The second is compact quasicompactness plus the annihilator theorem.  The
third is Lemma \ref{lem:r27-a2-dolgopyat}.  The fourth follows from
Lemmas \ref{lem:r27-a2-operator-uni} and \ref{lem:r27-a2-bad}; branch
derivatives contribute only $n^C$.  Integrability follows from $M>d_c$.
For the intermediate integral, split at
$|b|=\exp(n^{1/(2c_1)})$: below the split (A2.11) is superpolynomially small;
above it use the coarea estimate, which is valid in every frequency and whose
tail is $O(e^{-(M-d_c)n^{1/(2c_1)}})$.  This proves the stated local-theorem
budget.
\end{proof}

\subsection{Raw mixed local theorem and source insertions}

Let $Z_n=(S_n\kappa,S_nR,S_n\tau)$, with counting measure in the three lattice
coordinates and Lebesgue measure in the $d_c$ roof coordinates.

\begin{theorem}[Source-uniform mixed local limit theorem]
\label{thm:r27-a2-llt}
For strong-space insertions $f,g$ and central deviations
$y_n=O(\sqrt n)$,
\[
 \int f\,1_{\{Z_{n,\mathbb Z}=z_n\}}\,
 \delta_{Z_{n,\mathbb R}}(dt)\,g\circ F^n\,d\mu
 =
 {e^{-\frac1{2n}y_n^\mathsf T\Sigma^{-1}y_n}
  \over(2\pi n)^{(3+d_c)/2}\sqrt{\det\Sigma}}
 \bigl(\mu f\,\mu g+O(n^{-1/2})\bigr)\,dt.                \tag{A2.13}
\]
The estimate is uniform on the certified parameter family and under two
derivatives of a bounded finite-memory Feynman--Kac source.
\end{theorem}

\begin{proof}
Fourier inversion is split according to (A2.12).  In the
$n^{-1/2}$ ball, the Riesz projector and (A2.10) give the Gaussian with a
uniform Edgeworth remainder.  The other three regions are
$o(n^{-(3+d_c)/2})$ by Theorem \ref{thm:r27-a2-fourier}.  Source insertions
are multipliers in the fixed strong space; the same derivative budget and
Cauchy estimates give uniformity.  No smoothing of the roof variable is used.
\end{proof}

The proof now contains the missing anisotropic pairing, branch matching,
bad-word oscillation, exact arithmetic certificate, high differentiability
budget, and all-frequency integrable majorant.
